% =============================================================================
% Patch: two-sided whitening ledger for §4
%
% Insert immediately after Remark~\ref{rem:whitening-loss-scope} in
% §4 ``Finite-scale local blocks and whitening''.
%
% Purpose: Lemma~\ref{lem:whitening-loss} controls only upper-bound
% amplification by same-point whitening.  The following ledger records the
% lower-bound attenuation paid when a toy/source lower bound is first proved
% in unwhitened coordinates and then transferred into the whitened coordinates.
% =============================================================================

\subsection{Two-sided whitening ledger}
\label{subsec:two-sided-whitening-ledger}

The one-sided estimate of Lemma~\ref{lem:whitening-loss} prevents
same-point whitening from enlarging upper bounds.  Lower bounds require a
separate ledger, because left and right multiplication by
\(G_{m,\pm}^{-1/2}\) can attenuate a nonzero matrix.  Define the whitening
map on \(2\times2\) real matrices by
\begin{equation}
\label{eq:whitening-map}
        \mathsf W_{m,s}(X)
        :=G_{m,-}(s)^{-1/2}\,X\,G_{m,+}(s)^{-1/2}.
\end{equation}

\begin{lemma}[Spectral size of the finite same-point blocks]
\label{lem:finite-same-point-spectral-size}
For \(\Ph_T=\theta\), uniformly for \((m,s)\in\mathcal D_T\),
\begin{align}
\label{eq:G-lambda-min-asymptotic}
\lambda_{\min}\bigl(G_{m,\pm}(s)\bigr)
&=\frac{2q(t_\pm)}{\pi}\,(1+o_T(1)),\\
\label{eq:G-lambda-max-asymptotic}
\lambda_{\max}\bigl(G_{m,\pm}(s)\bigr)
&=\frac{q(t_\pm)^3}{6\pi}\,(1+o_T(1)).
\end{align}
Consequently, after increasing the global lower-height cutoff, there are
absolute constants \(c,C>0\) such that, uniformly on \(\mathcal D_T\),
\begin{equation}
\label{eq:G-spectral-comparable}
        c\,q(t_\pm)
        \le \lambda_{\min}\bigl(G_{m,\pm}(s)\bigr)
        \le C\,q(t_\pm),
        \qquad
        c\,q(t_\pm)^3
        \le \lambda_{\max}\bigl(G_{m,\pm}(s)\bigr)
        \le C\,q(t_\pm)^3.
\end{equation}
\end{lemma}

\begin{proof}
By Lemma~\ref{lem:finite-blocks-exact-interpretation},
\(G_{m,\pm}(s)=J(t_\pm)\).  Since \((m,s)\in\mathcal D_T\), both
\(t_\pm\in I_T\), and the theta derivative asymptotics of
Lemma~\ref{lem:theta-derivative-asymptotics} are uniform at these two
points.  The proof of Lemma~\ref{lem:same-point-gram-positivity} therefore
applies uniformly with \(T\) replaced by any \(t\in I_T\), giving
\eqref{eq:G-lambda-min-asymptotic}.

For the largest eigenvalue, use
\[
        \lambda_{\max}(J(t))
        =\Tr J(t)-\lambda_{\min}(J(t)).
\]
By \eqref{eq:J-trace-det}, uniformly for \(t\in I_T\),
\[
        \Tr J(t)
        =\frac{2q(t)^3+24q(t)+q''(t)}{12\pi}
        =\frac{q(t)^3}{6\pi}\,(1+o_T(1)),
\]
while \(\lambda_{\min}(J(t))=O(q(t))\).  Since \(q(t)\to\infty\)
uniformly on \(I_T\), the trace term of size \(q(t)^3\) dominates,
which gives \eqref{eq:G-lambda-max-asymptotic}.  The two-sided
comparability bounds \eqref{eq:G-spectral-comparable} follow by enlarging
the lower-height cutoff once more.
\end{proof}

\begin{lemma}[Two-sided whitening ledger]
\label{lem:two-sided-whitening-ledger}
For \(\Ph_T=\theta\), after increasing the global lower-height cutoff,
there are absolute constants \(c,C>0\) such that for every
\((m,s)\in\mathcal D_T\) and every real \(2\times2\) matrix \(X\),
\begin{equation}
\label{eq:two-sided-whitening-ledger}
        c\,\bigl(q(t_-)q(t_+)\bigr)^{-3/2}\,
        \|X\|_{\op}
        \le
        \|\mathsf W_{m,s}(X)\|_{\op}
        \le
        C\,\bigl(q(t_-)q(t_+)\bigr)^{-1/2}\,
        \|X\|_{\op}.
\end{equation}
In particular, the lower-bound attenuation of two-sided whitening is at
worst
\begin{equation}
\label{eq:two-sided-whitening-log-loss}
        \bigl(q(t_-)q(t_+)\bigr)^{-3/2}
        \asymp (\log T)^{-3}
\end{equation}
uniformly on \(\mathcal D_T\).
\end{lemma}

\begin{proof}
Let
\[
        A=G_{m,-}(s)^{-1/2},\qquad B=G_{m,+}(s)^{-1/2}.
\]
The upper bound follows from submultiplicativity:
\[
        \|AXB\|_{\op}
        \le \|A\|_{\op}\,\|X\|_{\op}\,\|B\|_{\op}
        = \lambda_{\min}(G_{m,-})^{-1/2}
          \lambda_{\min}(G_{m,+})^{-1/2}\,\|X\|_{\op}.
\]
Using the lower-eigenvalue estimate in
Lemma~\ref{lem:finite-same-point-spectral-size} gives the right-hand
inequality in \eqref{eq:two-sided-whitening-ledger}.

For the lower bound, write
\[
        X=A^{-1}(AXB)B^{-1}.
\]
Therefore
\[
        \|X\|_{\op}
        \le \|A^{-1}\|_{\op}\,\|AXB\|_{\op}\,\|B^{-1}\|_{\op}
        =\lambda_{\max}(G_{m,-})^{1/2}
         \lambda_{\max}(G_{m,+})^{1/2}\,\|AXB\|_{\op}.
\]
The upper-eigenvalue estimate in
Lemma~\ref{lem:finite-same-point-spectral-size} gives
\[
        \|AXB\|_{\op}
        \ge
        c\,\bigl(q(t_-)q(t_+)\bigr)^{-3/2}\,
        \|X\|_{\op},
\]
which is the left-hand inequality in
\eqref{eq:two-sided-whitening-ledger}.  Finally,
Lemma~\ref{lem:theta-derivative-asymptotics} gives
\(q(t_\pm)\asymp\log T\) uniformly on \(I_T\), and hence
\eqref{eq:two-sided-whitening-log-loss}.
\end{proof}

\begin{corollary}[Polynomial ledger form]
\label{cor:two-sided-whitening-polynomial-ledger}
Assume the packet scale has an upper phase-slope ledger
\begin{equation}
\label{eq:upper-phase-slope-ledger}
        q(t)\le Q(T)^{C_q^+}
        \qquad (t\in I_T)
\end{equation}
for some absolute exponent \(C_q^+\).  Then for all retained
\((m,s)\in\mathcal D_T\) and all real \(2\times2\) matrices \(X\),
\begin{equation}
\label{eq:two-sided-whitening-Q-ledger}
        \|\mathsf W_{m,s}(X)\|_{\op}
        \ge
        c\,Q(T)^{-3C_q^+}\,\|X\|_{\op},
\end{equation}
with an absolute constant \(c>0\).  Thus any pre-whitening lower bound
\(\|X\|_{\op}\ge Q(T)^{-A}\) becomes the post-whitening lower bound
\begin{equation}
\label{eq:pre-to-post-whitening-lower-bound}
        \|\mathsf W_{m,s}(X)\|_{\op}
        \ge c\,Q(T)^{-(A+3C_q^+)}.
\end{equation}
\end{corollary}

\begin{proof}
Combine the lower bound in Lemma~\ref{lem:two-sided-whitening-ledger} with
\eqref{eq:upper-phase-slope-ledger} at \(t_\pm\):
\[
        \bigl(q(t_-)q(t_+)\bigr)^{-3/2}
        \ge Q(T)^{-3C_q^+}.
\]
This gives \eqref{eq:two-sided-whitening-Q-ledger}; applying it to a
matrix satisfying \(\|X\|_{\op}\ge Q(T)^{-A}\) gives
\eqref{eq:pre-to-post-whitening-lower-bound}.
\end{proof}

\begin{remark}[When the ledger is consumed]
\label{rem:when-two-sided-whitening-ledger-is-consumed}
Lemma~\ref{lem:two-sided-whitening-ledger} is needed only when a toy,
source-side, or finite-algebraic lower bound is first proved in
unwhitened coordinates and then transported through
\(\mathsf W_{m,s}\).  If a visibility theorem is stated directly for the
already whitened block \(\widehat\Om_\z\), then no additional whitening
loss is charged.  The same caveat applies to any fixed finite-dimensional
norm other than \(\|\cdot\|_{\op}\): because all norms on \(M_2(\mathbb R)\)
are equivalent, the constants in the ledger change, but the powers of
\(q(t_\pm)\), \(\log T\), and \(Q(T)\) do not.
\end{remark}
